    % =========================================================
\subsection{City Planner / Local Authority}
% =====================================================================
For a city planner, two questions are most informative: \emph{which areas matter most to the
model}, and \emph{which built-environment features it relies on, and where}
(Figure~\ref{fig:stakeholder-planner}).\par\medskip
% City Planner stakeholder card (Chapter 6).
% Question: which areas matter most to the model, and where does it rely on built-environment
%           features, and which ones?
% Measure: share of each LSOA's total local Feature importance (|local coef| over all 30
% players) from the 12 built-environment players (Feature branch); map 01 adds signed
% GeoSensitivity (outputs/geo_area_data_eval/geo_coefficients_N4000.csv) for context.
%
% Text: generated by the dashboard LLM pipeline (generate_planner_card_narrative.py): master
% prompt geosmile-hull-v8 + card-mode addendum, the author-written card as an illustrative
% example. The model writes the headline, tiles and sections and marks figures ([[n:]] -> \pv)
% and key terms ([[t:]] -> \kt) itself. The generated headline (\cardHeadline) is not shown
% on the card. Prompts, evidence, raw output and checks:
% planner_card_llm.json. Hand-written version: planner_card_authored.tex
%
% Preamble: \usepackage[most]{tcolorbox}, xcolor, graphicx.
% Copy planner_geo_map.pdf and planner_built_map.pdf (plot_planner_maps.py) into Chapter6/figs/,
% and planner_card_llm_content_openai.tex into Chapter6/.
% =====================================================================

\definecolor{cardInk}{HTML}{12171C}
\definecolor{cardMuted}{HTML}{57646D}
\definecolor{cardTeal}{HTML}{1F4E5F}
\definecolor{cardPink}{HTML}{FCE4EC}
\definecolor{cardWine}{HTML}{8A0033}
\definecolor{cardRule}{HTML}{D9DEE2}

% Figures and key terms the model marked: one pink highlight in the running text's own font,
% size and weight; minimal padding so line spacing is unchanged. \colorbox does not break lines.
\providecommand{\hlpink}[1]{\begingroup\setlength{\fboxsep}{0.6pt}\colorbox{cardPink}{\textcolor{cardWine}{#1}}\endgroup}
\providecommand{\pv}[1]{\hlpink{#1}}
\providecommand{\kt}[1]{\hlpink{#1}}
\providecommand{\roletag}[1]{\tcbox[on line, colback=cardTeal, colframe=cardTeal, arc=2.5mm,
  boxsep=0pt, left=6pt, right=6pt, top=3pt, bottom=3pt]{\color{white}\sffamily\bfseries #1}}
\providecommand{\kicker}[1]{{\color{cardMuted}\sffamily\small\bfseries #1}}
% Stat tile: figure over a short label, four to a row, equal heights; body text size.
\providecommand{\tile}[2]{\begin{tcolorbox}[enhanced, nobeforeafter, width=\cardtilewidth,
  equal height group=\cardtilegroup, colback=cardPink, frame hidden, arc=2mm,
  left=5pt, right=3pt, top=3pt, bottom=3pt, valign=top, before upper=\raggedright]
  {\bfseries\color{cardWine}#1}\par\vspace{1pt}
  {\sffamily\color{cardMuted}#2}\end{tcolorbox}}

% Macros use \providecommand so several stakeholder cards can sit in one chapter;
% each card sets its own tile-height group and content file.
\def\cardtilegroup{plannertiles}
\def\cardtilewidth{0.2375\linewidth}
% Placeholders, replaced by the generated content file.
\providecommand{\cardHeadline}{[headline]}
\providecommand{\cardSecA}{[section 01]}
\providecommand{\cardSecB}{[section 02]}
\providecommand{\cardFooter}{[footer]}
\providecommand{\cardModel}{[model]}
\providecommand{\cardTileA}{\tile{--}{--}}
\providecommand{\cardTileB}{\tile{--}{--}}
\providecommand{\cardTileC}{\tile{--}{--}}
\providecommand{\cardTileD}{\tile{--}{--}}
\input{Chapter6/figs/planner_card_llm_content_openai.tex}

\begin{figure}[!htbp]
\begin{tcolorbox}[enhanced, colback=white, colframe=cardRule, arc=3mm, boxrule=0.6pt,
                  left=3mm, right=3mm, top=2.5mm, bottom=2.5mm, before skip=0pt, after skip=2pt,
                  fontupper=\sffamily\normalsize\color{cardInk}]
\setlength{\parskip}{0pt}\raggedright%
\roletag{CITY PLANNER}\par\smallskip
{\color{cardMuted}Important areas and built-environment reliance}\par\vspace{4pt}
\begin{minipage}[t]{0.485\linewidth}
  \kicker{01 \textperiodcentered{} Most important areas}\par\vspace{2pt}
  {\centering\includegraphics[width=0.88\linewidth]{Chapter6/figs/planner_geo_map.png}\par}\par\vspace{2pt}
  \cardSecA\par
\end{minipage}\hfill
\begin{minipage}[t]{0.485\linewidth}
  \kicker{02 \textperiodcentered{} Built-environment share}\par\vspace{2pt}
  {\centering\includegraphics[width=0.88\linewidth]{Chapter6/figs/planner_built_map.png}\par}\par\vspace{2pt}
  \cardSecB\par
\end{minipage}\par\vspace{3pt}
\kicker{Key figures}\par\vspace{2pt}
\noindent\cardTileA\hfill\cardTileB\hfill\cardTileC\hfill\cardTileD\par\vspace{4pt}
{\color{cardRule}\rule{\linewidth}{0.4pt}}\par\vspace{1pt}
{\color{cardMuted}\footnotesize \cardFooter{} Map 01: signed GeoSensitivity (sign is not a crime
direction). Map 02: share of local importance, not an effect on crime. Text and \pv{highlights}
generated by \cardModel{} (prompt geosmile-hull-v8, card mode).\par}
\end{tcolorbox}
\caption{City Planner card: GeoSensitivity (left) and built-environment share of local importance
(right); text generated by the LLM pipeline.}
\label{fig:stakeholder-planner}
\end{figure}
